\section*{Agradecimientos}
\label{sec:acknowledgement}
\addcontentsline{toc}{section}{Agradecimientos}

Quiero expresar mi más sincero agradecimiento a mi tutor, Carlos, por confiarme un tema tan estimulante y por acompañarlo después con criterio, paciencia y una disponibilidad poco habituales.
Su entusiasmo por el proyecto y la confianza que depositó en mí desde el principio fueron, en los momentos de duda, el empujón que necesitaba para seguir adelante.
A él le debo, además, la financiación del proyecto, a través del PIE consiguió una ayuda económica que, aunque modesta (no da para retirarse, pero algo es algo), tuvo el mérito de convertir muchas horas de entusiasmo en algo, al menos en parte, oficialmente remunerado.

A Rodri y Adri, compañeros de trinchera durante todo el máster.
Con vosotros las horas de trabajo pesaron la mitad y las de descanso valieron el doble: los cafés que se alargaban, las prácticas a última hora, las quejas compartidas y las risas que las desactivaban.
Buena parte de lo que me llevo de este año no está en ningún temario, sino en el equipo que hicimos dentro y fuera del aula.
Gracias por hacer del camino algo que apetecía recorrer, y mucha suerte y mucha fuerza en todo lo que venga.
Os llevo en el corazón para siempre.

Y, sobre todo, a mi familia y a mis amigos, a quienes dedico este trabajo: ellos sostuvieron todo lo demás (y a mí) mientras yo me ocupaba de estas páginas.
Lo que les debo no cabe en un agradecimiento, por eso les dedico el trabajo entero.

\vspace{\fill}
\begin{spacing}{1.0}
    Este trabajo se enmarca en el \gls{pie} IE26.6102~\cite{pie_upmmatch_2026}, \emph{UPM Match: Sistema de Recomendación de Tutores y Colaboraciones} (convocatoria 2025-2026, línea E1: Integración de la IA en las actividades docentes), de la \gls{upm}.
\end{spacing}
